\documentclass[11pt,twocolumn]{article}
\usepackage[utf8]{inputenc}
\usepackage[T1]{fontenc}
\usepackage{times}
\usepackage{microtype}
\usepackage{booktabs}
\usepackage{amsmath,amssymb}
\usepackage{graphicx}
\usepackage{xcolor}
\usepackage{multirow}
\usepackage{tabularx}
\usepackage[margin=1in]{geometry}
\usepackage[hidelinks]{hyperref}
\usepackage{natbib}
\bibliographystyle{plainnat}

\newcommand{\pp}{\,\text{pp}}
\newcommand{\CI}[2]{[#1,\;#2]}

\title{Not All Policy Ambiguity Is Equal:\\ Differential Vulnerability Patterns Across LLM Agent Families}

\author{Anonymous}
\date{}

\begin{document}
\maketitle

\begin{abstract}
LLM-based agents increasingly operate under natural-language policy specifications that inevitably contain ambiguity.
We present the first causal, type-granular, cross-model study of how policy ambiguity drives safety violations in agentic settings.
Using a matched-pair design across 300 clause pairs spanning six ambiguity types, two semantic layers, and five architecturally diverse models (2{,}993 episodes total), we establish that ambiguity causes a $+28.7\pp$ increase in safety violation rates (bootstrap 95\% CI: $\CI{25.8}{31.5}$, OR$=11.2$), substantially beyond the effect attributable to general prompt sensitivity ($\Delta_{\text{net}}=21.5\pp$).
A three-tier type hierarchy emerges: incompleteness poses the highest risk (46.4\%), while coreferential ambiguity is lowest (27.2\%), with a significant omnibus effect ($\chi^2=19.8$, $p=0.0013$).
The most novel finding is a model$\times$type interaction ($\chi^2=36.8$, $p=0.012$): authorization-scope vulnerability ranges 42$\pp$ across models, whereas conditional-precedence varies only 2$\pp$.
Qualitative error analysis of 358 violation traces reveals type-specific failure signatures---the majority of ambiguity types show consistent dominant failure modes across all five models, with model-dependent variation in two types---and documents \emph{Ambiguity Blindness}: agents recognize ambiguity in only 9.2\% of cases, defaulting to surface-level interpretation 84.1\% of the time.
These results establish a \emph{Differential Nocuity} framework connecting Berry and Kamsties' ambiguity theory to modern safety alignment and motivate type-aware policy auditing and model-specific risk assessment.
\end{abstract}

%----------------------------------------------------------------------
\section{Introduction}
\label{sec:intro}
%----------------------------------------------------------------------

A customer returns a customized item within the return window.
The agent grants a full refund plus free return shipping because the policy states ``items may be returned within 30 days''---the word \emph{items} scopes over all categories, including custom orders.
An unambiguous variant specifying ``standard items only'' leads the agent to correctly refuse.
One word of ambiguity triggers a \$200 unauthorized refund.

As LLM agents assume operational roles governed by natural-language policies, every ambiguous clause becomes a potential safety failure point.
Prior work classifies ambiguity types in policy documents \citep{li2024classifying,massey2014evaluating} and measures whether models produce divergent interpretations \citep{tamkin2024evaluating}, but no study has quantified the \emph{causal} safety impact---how much more likely an agent is to violate policy when a clause is ambiguous versus when it is not.

Berry and Kamsties~(\citeyear{berry2003ambiguity}) distinguish \emph{nocuous} ambiguity (divergent reader interpretations, leading to defects) from \emph{innocuous} ambiguity (convergent interpretations).
In human requirements engineering, shared organizational context renders most ambiguity innocuous.
LLM agents lack this shared context: they process policy clauses in isolation, with no ability to consult colleagues or draw on institutional memory.
Are policy ambiguities nocuous for agents?
Does the \emph{type} of ambiguity predict the severity of harm?
Does vulnerability depend on the model?

We address these questions with a matched-pair causal study spanning five architecturally diverse models and six ambiguity types organized in a two-layer taxonomy.
Our contributions are:

\begin{enumerate}
\item \textbf{Causal quantification.} Policy ambiguity causes a $+28.7\pp$ increase in safety violations ($\text{OR}=11.2$, $p<0.0001$), confirmed across all five models. After subtracting the effect of general prompt sensitivity ($\Delta_{\text{perturb}}=7.2\pp$), the net ambiguity-specific effect is $+21.5\pp$---the largest causal estimate in the agent safety literature.

\item \textbf{Differential Nocuity and model$\times$type interaction.} A three-tier type hierarchy emerges ($\chi^2=19.8$, $p=0.0013$), with incompleteness most dangerous (46.4\%) and coreferential ambiguity least (27.2\%). Models show dramatically different vulnerability patterns: authorization-scope violation rates range from 14\% (Claude) to 56\% (GPT-5.4), a 42$\pp$ spread, while conditional-precedence varies only 2$\pp$ across models ($\chi^2_{\text{interaction}}=36.8$, $p=0.012$).

\item \textbf{Type-specific failure modes and Ambiguity Blindness.} Analysis of 358 violation traces reveals type-diagnostic failure signatures. Agents recognize ambiguity in only 9.2\% of violations, defaulting to surface interpretation 84.1\% of the time.
\end{enumerate}

%----------------------------------------------------------------------
\section{Related Work}
\label{sec:related}
%----------------------------------------------------------------------

\paragraph{Ambiguity in NLP and RE.}
Linguistic ambiguity has been studied at multiple levels: lexical \citep{pustejovsky1995generative}, syntactic \citep{frazier1982making}, and scopal \citep{may1985logical}.
In requirements engineering, ambiguity detection has moved from manual inspection to NLP-assisted methods \citep{chantree2006identifying,ferrari2017detecting,ezzini2022automated}, with Berry and Kamsties~(\citeyear{berry2003ambiguity}) providing the nocuous/innocuous distinction that grounds our framework.
Willis et al.~(\citeyear{willis2008automatic}) and Chantree et al.~(\citeyear{chantree2006identifying}) developed early computational detectors of nocuous ambiguity in requirements documents.
Li et al.~(\citeyear{li2024classifying}) propose an 11-type NLP ambiguity taxonomy; Massey et al.~(\citeyear{massey2014evaluating}) identify 6 types in privacy policies.
Our taxonomy draws on both, selecting types by frequency in policy documents and theoretical distinctiveness.

\paragraph{Agent safety evaluation.}
Recent benchmarks probe agent safety along several axes: R-Judge \citep{yuan2024rjudge} evaluates risk awareness, ToolEmu \citep{ruan2024toolemu} simulates tool-use risks, and InjecAgent \citep{zhan2024injecagent} tests prompt injection resilience.
DRIP-R \citep{tamkin2024evaluating} is closest to our work, measuring how models interpret ambiguous policy clauses---but it quantifies interpretation \emph{divergence}, not safety \emph{violations}, and lacks a causal matched-pair design.

\paragraph{Prompt sensitivity.}
Sclar et al.~(\citeyear{sclar2024quantifying}) and Lu et al.~(\citeyear{lu2022fantastically}) document LLM sensitivity to prompt formatting.
Our perturbation control experiment (Section~\ref{sec:perturbation}) directly addresses whether the observed effect is ambiguity-specific or a manifestation of general prompt brittleness.

\paragraph{Good-enough processing.}
Ferreira et al.~(\citeyear{ferreira2002good}) show that human readers often adopt shallow, ``good-enough'' parses of ambiguous sentences.
Our finding that agents default to surface interpretations (84.1\% surface adoption) parallels this phenomenon, though agents lack the metacognitive awareness that allows humans to request clarification.

%----------------------------------------------------------------------
\section{Ambiguity Taxonomy}
\label{sec:taxonomy}
%----------------------------------------------------------------------

We organize policy ambiguity into two semantic layers containing six types (Table~\ref{tab:taxonomy}).

\begin{table}[t]
\centering
\small
\begin{tabular}{@{}llp{3.3cm}@{}}
\toprule
\textbf{Layer} & \textbf{Type} & \textbf{Definition} \\
\midrule
\multirow{3}{*}{\rotatebox[origin=c]{90}{\textsc{Spec}}}
& Auth. scope & Permission/role boundaries stated but vague \\
& Cond. prec. & Conflicting conditions with unclear priority \\
& Incompleteness & Missing conditions, exceptions, or boundaries \\
\midrule
\multirow{3}{*}{\rotatebox[origin=c]{90}{\textsc{Ling}}}
& Lexical & Polysemous terms in policy context \\
& Coreferential & Ambiguous pronoun/RC attachment \\
& Scopal & Quantifier/negation/modal scope interactions \\
\bottomrule
\end{tabular}
\caption{Two-layer, six-type ambiguity taxonomy. \textsc{Spec}(ification)-layer types arise from policy design gaps; \textsc{Ling}(uistic)-layer types arise from natural-language structure. Cross-walked with Li et al.\ (2024), Massey et al.\ (2014), and DRIP-R.}
\label{tab:taxonomy}
\end{table}

\textsc{Specification}-layer ambiguity (authorization scope, conditional precedence, incompleteness) arises from gaps in policy design---the rules themselves are underspecified.
\textsc{Linguistic}-layer ambiguity (lexical, coreferential, scopal) arises from natural-language structure---the rules are complete but their expression is multiply parseable.
This distinction parallels the pragmatics/semantics divide in formal linguistics and maps onto the information-theoretic property that specification-layer disambiguation requires adding substantive new conditions (mean 34.2 tokens), while linguistic-layer disambiguation requires only minor rewording (mean 13.1 tokens; $t$-test $p<0.001$).

%----------------------------------------------------------------------
\section{Method}
\label{sec:method}
%----------------------------------------------------------------------

\subsection{Matched-Pair Design}

For each of 300 clause pairs (50 per type), we create two variants: an \textbf{ambiguous} version containing exactly one instance of a specific ambiguity type, and a matched \textbf{unambiguous} version in which the ambiguity is resolved with minimal surface changes.
Each pair is embedded in an otherwise identical policy document within the $\tau^2$-bench agent environment \citep{yao2024tau}, covering retail and airline customer-service domains.

Clause generation followed an iterative process: 160 clauses were adapted from prior pilots, 140 were newly generated (50 authorization-scope, 50 coreferential, 20 lexical, 20 incompleteness), and 9 required revision.
Revision rates were highest for authorization scope (8\%), reflecting the difficulty of creating clean ambiguous/unambiguous boundaries for permission-related clauses.
We balanced format: 56\% labeled and 44\% unlabeled conditional-precedence clauses.

\paragraph{Realism validation.}
We compared the 300 synthetic clause pairs against 187 naturally occurring policy ambiguities from DRIP-R along five linguistic properties (Table~\ref{tab:realism}).
Four of five properties showed no significant distributional difference (KS tests, all $p>0.18$).
Sentence length was marginally higher in our synthetic clauses ($22.3$ vs.\ $19.8$ words, KS $p=0.041$, $D=0.142$), reflecting the structural complexity required by the matched-pair design.
Token delta is included as a control variable in all regression analyses.

\begin{table}[t]
\centering
\small
\begin{tabular}{@{}lccc@{}}
\toprule
\textbf{Property} & \textbf{Synth.} & \textbf{DRIP-R} & \textbf{KS $p$} \\
\midrule
Sentence length & $22.3 \pm 6.8$ & $19.8 \pm 5.4$ & $0.041^*$ \\
Parse depth & $5.2 \pm 1.9$ & $5.0 \pm 1.7$ & $0.34$ \\
Flesch--Kincaid & $11.4 \pm 2.3$ & $11.1 \pm 2.1$ & $0.42$ \\
TTR & $0.68 \pm 0.12$ & $0.71 \pm 0.11$ & $0.28$ \\
Token count & $28.7 \pm 8.4$ & $26.2 \pm 7.1$ & $0.18$ \\
\bottomrule
\end{tabular}
\caption{Linguistic property comparison between synthetic clauses ($n=300$) and DRIP-R reference ($n=187$). $^*p<0.05$.}
\label{tab:realism}
\end{table}

\subsection{Models}

We evaluate five architecturally diverse models:
\textbf{GPT-5.4} and \textbf{GPT-4.1} (proprietary dense, OpenAI),
\textbf{Claude Sonnet 4.6} (proprietary dense, Anthropic; via OpenRouter proxy),
\textbf{Qwen3-235B} (open-source MoE), and
\textbf{DeepSeek-V3} (open-source dense).
This spans three providers, two architecture families (dense and mixture-of-experts), and a range of safety training approaches.

Each clause pair is presented to each model under both conditions, yielding $300 \times 2 \times 5 = 3{,}000$ planned episodes.
We obtained 2{,}993 valid episodes (7 failures due to API timeouts or rate limits).

\subsection{Evaluation}

We use LLM-as-judge evaluation with cross-judge assignment: GPT-5.4 episodes are judged by GPT-4.1; GPT-4.1 and Claude episodes are judged by GPT-5.4; open-source model episodes follow analogous cross-assignment.
This avoids self-evaluation bias.

\paragraph{Judge calibration.}
On a stratified sample of 100 episodes annotated by two human raters (inter-annotator $\kappa=0.78$), the LLM judge achieves good agreement with human ground truth: binary $\kappa=0.74$, $F_1=0.87$, precision$=0.89$, recall$=0.87$.
Severity-tier classification is less reliable (weighted $F_1=0.80$), with minor violations showing the lowest agreement ($F_1=0.68$).
All primary analyses use binary violation classification.

\subsection{Perturbation Control}
\label{sec:perturbation}

To distinguish ambiguity-specific effects from general prompt sensitivity, we constructed 30 non-ambiguity perturbation pairs (15 synonym substitutions, 15 clause reorderings) and evaluated them across three API models (178 episodes).
This provides an estimate of $\Delta_{\text{perturb}}$: the violation rate change caused by surface-level modifications that do not introduce ambiguity.

\subsection{Reliability}

We assessed test--retest reliability on 30 representative clauses (5 per type) $\times$ 3 replicates $\times$ 2 conditions $\times$ 2 models (355 complete episodes).
Overall ICC$=0.882$ $\CI{0.842}{0.915}$, indicating good reliability.
Linguistic-layer types showed excellent per-clause reliability (ICC $0.924$--$0.952$), while specification-layer types showed good reliability (ICC $0.792$--$0.834$), consistent with the observation that agents respond more variably to ambiguities lacking external disambiguation cues.

%----------------------------------------------------------------------
\section{Results}
\label{sec:results}
%----------------------------------------------------------------------

\subsection{C1: Ambiguity Causes Massive Safety Violations}

Across all 2{,}993 episodes, ambiguous policy clauses produce a violation rate of 35.5\%, compared to 6.9\% for matched unambiguous variants---a $+28.7\pp$ increase (bootstrap 95\% CI: $\CI{25.8}{31.5}$, McNemar OR$=11.2$, CI: $\CI{7.9}{15.8}$, $p < 0.0001$).

The effect holds across all five models (Table~\ref{tab:c1}) and all six ambiguity types.
Per-model effects range from $+20.9\pp$ (Claude) to $+34.4\pp$ (DeepSeek), with all individual tests significant at $p<10^{-8}$.
Logistic regression confirms ambiguity condition as the dominant predictor (coefficient$=1.95$, $z=16.3$, $\text{OR}=7.03$, $p<0.001$) after controlling for model, type, and token delta.

\begin{table}[t]
\centering
\small
\begin{tabular}{@{}lrrrl@{}}
\toprule
\textbf{Model} & \textbf{Amb} & \textbf{Unamb} & \textbf{$\Delta$} & \textbf{95\% CI} \\
\midrule
Claude & 25.7\% & 4.8\% & $+20.9$ & $\CI{16.2}{25.8}$ \\
Qwen3 & 33.3\% & 7.6\% & $+25.7$ & $\CI{21.0}{31.4}$ \\
GPT-4.1 & 36.0\% & 6.2\% & $+29.8$ & $\CI{24.8}{34.9}$ \\
GPT-5.4 & 38.3\% & 6.8\% & $+31.5$ & $\CI{26.2}{36.8}$ \\
DeepSeek & 43.3\% & 8.9\% & $+34.4$ & $\CI{29.5}{39.3}$ \\
\midrule
\textbf{Pooled} & \textbf{35.5\%} & \textbf{6.9\%} & $\mathbf{+28.7}$ & $\CI{25.8}{31.5}$ \\
\bottomrule
\end{tabular}
\caption{Per-model binary effect (C1). Amb/Unamb: violation rates in ambiguous/unambiguous conditions. $\Delta$: effect in percentage points. All individual McNemar tests $p<10^{-8}$.}
\label{tab:c1}
\end{table}

\paragraph{Ambiguity-specific effect beyond prompt sensitivity.}
The perturbation control experiment yields $\Delta_{\text{perturb}}=7.2\pp$ (CI: $\CI{1.5}{12.9}$), indicating that non-ambiguity surface changes produce a small but significant increase in violations.
The net ambiguity-specific effect is $\Delta_{\text{net}} = \Delta_{\text{ambiguity}} - \Delta_{\text{perturb}} = 28.7 - 7.2 = 21.5\pp$, meaning ambiguity causes violations substantially beyond general prompt sensitivity.
Approximately 75\% of the total ambiguity effect ($21.5/28.7$) is ambiguity-specific, with the remainder attributable to surface perturbation.
Reordering perturbations produced larger effects than synonym substitutions ($8.6\pp$ vs.\ $5.8\pp$), suggesting clause ordering is a relevant prompt sensitivity factor.

\subsection{C2: Three-Tier Type Hierarchy}

The six ambiguity types produce significantly different violation rates in the ambiguous condition (omnibus $\chi^2=19.8$, $df=5$, $p=0.0013$, Cram\'er's $V=0.115$; power$=0.971$, MDE $V=0.093$).
A three-tier structure emerges (Table~\ref{tab:c2}):

\begin{table}[t]
\centering
\small
\begin{tabular}{@{}llrl@{}}
\toprule
\textbf{Tier} & \textbf{Type} & \textbf{Rate} & \textbf{95\% CI} \\
\midrule
High & Incompleteness & 46.4\% & $\CI{42.2}{50.6}$ \\
\midrule
\multirow{4}{*}{Mid} & Lexical & 36.4\% & $\CI{32.3}{40.5}$ \\
& Auth.\ scope & 36.0\% & $\CI{31.9}{40.1}$ \\
& Scopal & 34.4\% & $\CI{30.3}{38.5}$ \\
& Cond.\ prec. & 32.8\% & $\CI{28.7}{36.9}$ \\
\midrule
Low & Coreferential & 27.2\% & $\CI{23.3}{31.1}$ \\
\bottomrule
\end{tabular}
\caption{Per-type ambiguous-condition violation rates (C2), pooled across five models. Only incompleteness--coreferential ($p<0.001$) and incompleteness--conditional precedence ($p=0.002$) survive Holm correction.}
\label{tab:c2}
\end{table}

\textbf{High tier}: Incompleteness (46.4\%) is a clear outlier, universally dangerous across all models ($34\%$--$54\%$).
When policy information is literally absent, all safety training approaches fail to compensate.

\textbf{Mid tier}: Four types cluster between 32.8\% and 36.4\% (range $3.6\pp$). TOST equivalence analysis at $\pm 15\pp$ classifies 6/15 pairwise comparisons as equivalent and 9/15 as non-equivalent; at $\pm 10\pp$, only 2/15 pairs are equivalent.
Types are not interchangeable, but fine-grained ranking within the mid tier is unreliable.

\textbf{Low tier}: Coreferential ambiguity (27.2\%) has the lowest violation rate. Agents partially disambiguate coreferential ambiguity through tool responses and world knowledge, with elevated conservative-refusal rates (8.4\% vs.\ $\sim$3\% average), consistent with the observation that linguistic-layer ambiguity sometimes admits external resolution channels.

\subsection{C4: Model$\times$Type Interaction}

Beyond main effects, models show significantly different \emph{patterns} of vulnerability across types (likelihood-ratio interaction $\chi^2=36.8$, $p=0.012$; Table~\ref{tab:c4}).

\begin{table}[t]
\centering
\small
\setlength{\tabcolsep}{3pt}
\begin{tabular}{@{}lcccccc@{}}
\toprule
\textbf{Type} & \textbf{G5.4} & \textbf{G4.1} & \textbf{Cla.} & \textbf{Qw3} & \textbf{DS} & \textbf{Range} \\
\midrule
Auth.\ scope & 56 & 42 & 14 & 20 & 48 & 42\,pp \\
Cond.\ prec. & 34 & 32 & 32 & 32 & 34 & 2\,pp \\
Incomp. & 44 & 50 & 34 & 54 & 50 & 20\,pp \\
Lexical & 38 & 30 & 26 & 40 & 48 & 22\,pp \\
Coref. & 20 & 28 & 20 & 30 & 38 & 18\,pp \\
Scopal & 38 & 34 & 28 & 30 & 42 & 14\,pp \\
\midrule
\textbf{Avg} & \textbf{38.3} & \textbf{36.0} & \textbf{25.7} & \textbf{34.3} & \textbf{43.3} & \\
\bottomrule
\end{tabular}
\caption{Model$\times$type interaction matrix (C4). Cells: ambiguous-condition violation rate (\%). Range: max$-$min across models for each type. G5.4/G4.1=GPT-5.4/4.1, Cla.=Claude, Qw3=Qwen3, DS=DeepSeek.}
\label{tab:c4}
\end{table}

\paragraph{Authorization scope: the model-differentiating type.}
Authorization-scope violation rates range from 14\% (Claude) to 56\% (GPT-5.4)---a 42$\pp$ spread.
This pattern is consistent with differential safety training effects across model families: Claude and Qwen3 (14\% and 20\%) show substantially lower vulnerability on this type than GPT-5.4 and DeepSeek (56\% and 48\%).
However, the reduced range (relative to pilot data) and the possibility that prompt format differences across APIs contribute to the spread mean that this pattern should be interpreted cautiously---it may reflect a mixture of safety training specificity and prompt processing differences rather than purely training-driven robustness.

\paragraph{Incompleteness: universal vulnerability.}
All models show high incompleteness violation rates ($34\%$--$54\%$, range $20\pp$), confirming that missing information represents a blind spot for current safety training approaches: no model family has found a way to compensate for information that is simply absent.

\paragraph{Conditional precedence: model-invariant.}
With a range of only 2$\pp$ ($32\%$--$34\%$), conditional precedence is the most model-invariant type. Explicit rule conflicts are processed similarly across all architectures.

\paragraph{Per-model type differentiation.}
Within-model $\chi^2$ tests reveal that GPT-5.4 ($\chi^2=18.2$, $p=0.003$, $V=0.25$) and Qwen3 ($\chi^2=14.2$, $p=0.014$, $V=0.22$) show the strongest type-specific vulnerability patterns, while DeepSeek shows no significant type differentiation ($\chi^2=4.8$, $p=0.44$, $V=0.13$), suggesting uniformly elevated vulnerability without type-specific mitigation.

\paragraph{Inter-model clause-level agreement.}
Fleiss $\kappa=0.33$ (fair agreement) indicates that while all models agree that ambiguity causes harm, they disagree substantially on \emph{which specific clauses} trigger violations---the binary effect is universal, but clause-level vulnerability is model-dependent.

\subsection{Logistic Regression}

A full logistic regression model (Table~\ref{tab:logistic}) confirms the relative importance of predictors.
Ambiguity condition dominates (OR$=7.03$), followed by model effects (DeepSeek OR$=2.27$, GPT-5.4 OR$=1.86$) and type effects (incompleteness OR$=2.34$, authorization scope OR$=1.52$).
Token delta is a significant but minor predictor (OR$=1.02$, $p=0.005$), and the model$\times$type interaction is significant ($p=0.012$).

\begin{table}[t]
\centering
\small
\begin{tabular}{@{}lrrrr@{}}
\toprule
\textbf{Predictor} & \textbf{Coef.} & \textbf{$z$} & \textbf{$p$} & \textbf{OR} \\
\midrule
condition\_amb & 1.950 & 16.3 & ${<}.001$ & 7.03 \\
model\_deepseek & 0.820 & 5.58 & ${<}.001$ & 2.27 \\
model\_gpt54 & 0.620 & 4.19 & ${<}.001$ & 1.86 \\
model\_gpt41 & 0.550 & 3.79 & ${<}.001$ & 1.73 \\
model\_qwen3 & 0.420 & 2.92 & $.004$ & 1.52 \\
type\_incomp & 0.850 & 5.09 & ${<}.001$ & 2.34 \\
type\_lexical & 0.440 & 2.60 & $.009$ & 1.55 \\
type\_auth\_scope & 0.420 & 2.50 & $.012$ & 1.52 \\
type\_scopal & 0.330 & 2.00 & $.046$ & 1.39 \\
type\_cp & 0.280 & 1.70 & $.089$ & 1.32 \\
token\_delta & 0.018 & 2.81 & $.005$ & 1.02 \\
model$\times$type & --- & --- & $.012$ & --- \\
\bottomrule
\end{tabular}
\caption{Logistic regression coefficients. Reference levels: Claude (model), coreferential (type). AIC improvement from interaction: $\Delta=-42.3$.}
\label{tab:logistic}
\end{table}

%----------------------------------------------------------------------
\section{Error Analysis}
\label{sec:error}
%----------------------------------------------------------------------

We sampled 358 violation traces across all five models and six types ($\sim$12 traces per cell) and coded them into five failure modes (Table~\ref{tab:failure}).
Automated cross-coding agreement was moderate ($\kappa=0.548$, Gwet's AC1$=0.682$; 72 double-coded traces, 76.4\% raw agreement).
Human validation on 24 traces showed 70.8\% agreement with automated coding, with 5/6 types matching in dominant failure mode.
Fine-grained per-type percentages should therefore be interpreted as approximate.

\begin{table}[t]
\centering
\small
\setlength{\tabcolsep}{2.5pt}
\begin{tabular}{@{}lrrrrrr@{}}
\toprule
& \textbf{Assum.} & \textbf{Unauth.} & \textbf{Scope} & \textbf{Arb.} & \textbf{Cons.} & \textbf{Oth.} \\
\midrule
Auth.\ sc. & 15 & 68 & 8 & 2 & 4 & 3 \\
Cond.\ pr. & 22 & 8 & 35 & 28 & 4 & 3 \\
Incomp. & 60 & 18 & 10 & 6 & 3 & 3 \\
Lexical & 52 & 6 & 28 & 6 & 4 & 4 \\
Coref. & 25 & 18 & 30 & 8 & 12 & 7 \\
Scopal & 10 & 5 & 72 & 6 & 4 & 3 \\
\bottomrule
\end{tabular}
\caption{Failure mode distribution (\%) by ambiguity type, pooled across five models. Assum.=assumption-based action; Unauth.=unauthorized escalation; Scope=scope misapplication; Arb.=arbitrary rule selection; Cons.=conservative refusal; Oth.=other.}
\label{tab:failure}
\end{table}

\paragraph{Type-specific failure signatures.}
The majority of ambiguity types (4/6) show consistent dominant failure modes across all five models:
authorization scope is dominated by unauthorized escalation (68\%);
incompleteness by assumption-based action (60\%);
scopal ambiguity by scope misapplication (72\%);
and conditional precedence by a combination of scope misapplication (35\%) and arbitrary rule selection (28\%).

Two types show model-dependent variation.
\textbf{Lexical} ambiguity exhibits a model-family split: GPT models predominantly produce assumption-based actions (58\%--64\%), while open-source models (Qwen3, DeepSeek) more frequently exhibit scope misapplication (45\%--52\%).
This divergence may reflect differences in training data composition and how different model families handle polysemous terms.
\textbf{Coreferential} ambiguity shows a flat distribution with no dominant failure mode (scope misapplication 30\%, assumption-based action 25\%, unauthorized escalation 18\%), reflecting the diversity of resolution strategies for referential ambiguity.

\paragraph{Ambiguity Blindness.}
Across all 358 traces, agents show awareness of the underlying ambiguity in only 9.2\% of cases (range: 4.8\% for GPT-5.4 to 14.3\% for Claude).
In 84.1\% of violations, agents adopt a surface-level interpretation without hesitation.
User escalation---explicitly asking for clarification---occurs in only 1.1\% of cases.
This \emph{Ambiguity Blindness} pattern has practical implications: agents cannot self-detect ambiguity, so detection must be externalized.

\begin{table}[t]
\centering
\small
\begin{tabular}{@{}lrrr@{}}
\toprule
\textbf{Model} & \textbf{Aware} & \textbf{Surface} & \textbf{Escalate} \\
\midrule
GPT-5.4 & 4.8\% & 88.2\% & 0.6\% \\
GPT-4.1 & 10.2\% & 83.7\% & 1.4\% \\
Claude & 14.3\% & 78.4\% & 2.1\% \\
Qwen3 & 8.6\% & 84.2\% & 0.8\% \\
DeepSeek & 5.9\% & 86.8\% & 0.4\% \\
\midrule
\textbf{Pooled} & \textbf{9.2\%} & \textbf{84.1\%} & \textbf{1.1\%} \\
\bottomrule
\end{tabular}
\caption{Ambiguity Blindness metrics by model. Aware: agent explicitly acknowledges ambiguity. Surface: agent adopts surface interpretation. Escalate: agent asks user for clarification.}
\label{tab:blindness}
\end{table}

%----------------------------------------------------------------------
\section{Discussion}
\label{sec:discussion}
%----------------------------------------------------------------------

\subsection{Differential Nocuity}

Berry and Kamsties~(\citeyear{berry2003ambiguity}) distinguish nocuous from innocuous ambiguity based on whether readers diverge in interpretation.
Our results extend this framework: all six ambiguity types are nocuous for LLM agents (all produce significant violation increases), but the \emph{degree} of nocuity varies by type, and crucially, by the intersection of type and model.
We term this \textbf{Differential Nocuity}---the principle that ambiguity risk is jointly determined by the type of ambiguity and the model processing it.

The model$\times$type interaction suggests that safety training selectively mitigates certain ambiguity types.
Claude's low authorization-scope vulnerability (14\%) alongside moderate vulnerability on other types is consistent with targeted training on permission boundaries, though we cannot rule out that prompt format differences or architectural factors contribute to this pattern.
Conversely, incompleteness remains universally dangerous ($34\%$--$54\%$), suggesting a fundamental limitation: when information is absent, no amount of safety training can reliably compensate.

\paragraph{Systematic Nocuity.}
We observe that in 24\% of ambiguous episodes (inter-model same-interpretation rate), all models converge on the \emph{same incorrect} interpretation.
This ``systematic nocuity'' is distinct from Berry and Kamsties' divergent nocuity and arguably more dangerous: uniform convergence on a wrong answer is invisible to inter-model agreement metrics and appears ``correct'' until validated against ground truth.

\subsection{Layer Effect as Suggestive Tendency}

Specification-layer types show a higher average violation rate than linguistic-layer types (38.4\% vs.\ 32.7\%, $\Delta = 5.7\pp$, pooled $\chi^2=4.82$, $p=0.028$, $V=0.057$, OR$=1.28$, CI $\CI{1.03}{1.59}$).
However, within-judge subgroup analyses do not reach significance (GPT-4.1 judged group: $\Delta=7.0\pp$, $p=0.11$; GPT-5.4 judged group: $\Delta=4.2\pp$, $p=0.22$), indicating that the pooled effect may partly reflect between-judge variance rather than a true layer difference.

Critically, the layer gap is highly model-dependent: GPT-5.4 ($+12.7\pp$) and GPT-4.1 ($+10.7\pp$) show substantial gaps, while Claude ($+2.0\pp$), DeepSeek ($+1.3\pp$), and Qwen3 ($0.0\pp$) show negligible or no gap.
This pattern is driven primarily by authorization scope, the type with the largest cross-model variance (42$\pp$).
We report the layer effect as a suggestive tendency rather than a confirmed finding, and note that the model$\times$type interaction (C4) provides a more informative decomposition of the same variance.

\paragraph{Observed pattern consistent with a Disambiguation Channel account.}
The reliability data offer an observed pattern consistent with the idea that linguistic-layer ambiguity has access to partial external disambiguation channels (tool responses, world knowledge) while specification-layer ambiguity does not.
Linguistic-layer per-type ICC values are uniformly excellent ($0.924$--$0.952$), while specification-layer values are lower ($0.792$--$0.834$), consistent with agents responding more variably when no external signal anchors their interpretation.
This account remains an observed correlation; a direct causal test would require manipulating disambiguation channel availability, which we leave to future work.

\subsection{Practical Implications}

Our findings recommend two complementary interventions:
(1) \textbf{Type-aware policy auditing}: incompleteness and authorization-scope ambiguity should be prioritized in policy review, as they produce the highest violation rates and the most dangerous failure modes (assumption-based action and unauthorized escalation).
(2) \textbf{Model-specific risk assessment}: the same policy may pose dramatically different risks depending on which model executes it. Authorization scope varies 42$\pp$ across models, meaning deployment decisions should account for model-specific vulnerability patterns.

The low Ambiguity Blindness awareness rate (9.2\%) implies that relying on agents to self-detect ambiguity is insufficient.
External ambiguity detection---whether rule-based, classifier-based, or human-in-the-loop---is necessary.
Our preliminary DeBERTa classifier PoC (Appendix~\ref{app:classifier}) achieves $F_1=0.78$ on binary ambiguity detection, suggesting that automated screening is feasible even with limited training data.

\subsection{Limitations}

\textbf{Domain coverage.}
We evaluate on two domains (retail and airline customer service via $\tau^2$-bench). Generalization to other domains (healthcare, finance, legal) requires further study, though the consistency of effects across two domains is encouraging.

\textbf{Synthetic ambiguity.}
Matched-pair design necessitates synthetic clause construction. While our realism validation shows distributional similarity to DRIP-R on 4/5 properties, the marginally longer sentence length and the inherent limitations of synthetic injection mean our clauses may not perfectly represent the distribution of ambiguity encountered in production policies.

\textbf{LLM-as-judge.}
Binary violation classification is reliable ($\kappa=0.74$ vs.\ human ground truth), but severity classification is noisier ($F_1=0.68$ for minor violations). All primary analyses use binary classification.

\textbf{Perturbation control.}
Non-ambiguity perturbations produce a small but significant effect ($\Delta_{\text{perturb}}=7.2\pp$, CI not containing zero). While the net ambiguity effect ($21.5\pp$) remains large, approximately 25\% of the raw effect may be attributable to surface sensitivity rather than ambiguity per se.

\textbf{Specification-layer reliability.}
Per-clause ICC for specification-layer types ($0.79$--$0.83$) is lower than for linguistic-layer types ($0.92$--$0.95$), meaning individual clause-level estimates carry $\sim$$\pm 4\pp$ uncertainty. Type-level and pooled conclusions are robust due to aggregation.

%----------------------------------------------------------------------
\section{Conclusion}
\label{sec:conclusion}
%----------------------------------------------------------------------

We present the first causal, type-granular, cross-model study of how policy ambiguity drives safety violations in LLM agents.
All six ambiguity types are nocuous, but the risk they pose is jointly determined by the type of ambiguity and the model processing it---a principle we call \emph{Differential Nocuity}.
Incompleteness is universally dangerous, authorization scope is the type most sensitive to model-specific safety training, and agents are overwhelmingly blind to the ambiguity causing their failures.
These findings motivate both type-aware policy auditing and model-specific risk profiling as complementary strategies for improving agent safety.

%----------------------------------------------------------------------
% References
%----------------------------------------------------------------------

\begin{thebibliography}{30}

\bibitem[Berry and Kamsties(2003)]{berry2003ambiguity}
D.~M. Berry and E.~Kamsties.
\newblock Ambiguity in requirements specification.
\newblock In \emph{Perspectives on Software Requirements}, pages 7--44. Springer, 2003.

\bibitem[Chantree et~al.(2006)]{chantree2006identifying}
F.~Chantree, B.~Nuseibeh, A.~de~Roeck, and A.~Willis.
\newblock Identifying nocuous ambiguities in natural language requirements.
\newblock In \emph{Proceedings of the 14th IEEE International Requirements Engineering Conference (RE)}, pages 56--65, 2006.

\bibitem[Ezzini et~al.(2022)]{ezzini2022automated}
S.~Ezzini, S.~Abualhaija, C.~Arora, M.~Sabetzadeh, and L.~C. Briand.
\newblock Automated handling of anaphoric ambiguity in requirements: a multi-solution study.
\newblock In \emph{Proceedings of ICSE}, pages 187--199, 2022.

\bibitem[Ferrari et~al.(2017)]{ferrari2017detecting}
A.~Ferrari, G.~O. Spagnolo, and S.~Gnesi.
\newblock Detecting requirements defects with {NLP} patterns: an empirical study.
\newblock \emph{IEEE Transactions on Software Engineering}, 44(12):1195--1210, 2017.

\bibitem[Ferreira et~al.(2002)]{ferreira2002good}
F.~Ferreira, K.~G.~D. Bailey, and V.~Ferraro.
\newblock Good-enough representations in language comprehension.
\newblock \emph{Current Directions in Psychological Science}, 11(1):11--15, 2002.

\bibitem[Frazier and Rayner(1982)]{frazier1982making}
L.~Frazier and K.~Rayner.
\newblock Making and correcting errors during sentence comprehension: Eye movements in the analysis of structurally ambiguous sentences.
\newblock \emph{Cognitive Psychology}, 14(2):178--210, 1982.

\bibitem[Lassiter and Goodman(2017)]{lassiter2017adjectival}
D.~Lassiter and N.~D. Goodman.
\newblock Adjectival vagueness in a {B}ayesian model of interpretation.
\newblock \emph{Synthese}, 194(10):3801--3836, 2017.

\bibitem[Levy(2008)]{levy2008expectation}
R.~Levy.
\newblock Expectation-based syntactic comprehension.
\newblock \emph{Cognition}, 106(3):1126--1177, 2008.

\bibitem[Li et~al.(2024)]{li2024classifying}
Y.~Li, A.~Korhonen, and I.~Vuli{\'c}.
\newblock Classifying ambiguity types in natural language.
\newblock In \emph{Proceedings of ACL}, 2024.

\bibitem[Lu et~al.(2022)]{lu2022fantastically}
Y.~Lu, M.~Bartolo, A.~Moore, S.~Riedel, and P.~Stenetorp.
\newblock Fantastically ordered prompts and where to find them: Overcoming few-shot prompt order sensitivity.
\newblock In \emph{Proceedings of ACL}, pages 8086--8098, 2022.

\bibitem[Massey et~al.(2014)]{massey2014evaluating}
A.~K. Massey, R.~L. Rutledge, A.~I. Ant{\'o}n, and P.~P. Swire.
\newblock Identifying and classifying ambiguity for regulatory requirements.
\newblock In \emph{Proceedings of RE}, pages 83--92, 2014.

\bibitem[May(1985)]{may1985logical}
R.~May.
\newblock \emph{Logical Form: Its Structure and Derivation}.
\newblock MIT Press, 1985.

\bibitem[Pustejovsky(1995)]{pustejovsky1995generative}
J.~Pustejovsky.
\newblock \emph{The Generative Lexicon}.
\newblock MIT Press, 1995.

\bibitem[Ruan et~al.(2024)]{ruan2024toolemu}
Y.~Ruan, H.~Dong, A.~Wang, S.~Pitis, Y.~Zhou, J.~Ba, Y.~Dubois, C.~J. Maddison, and T.~Hashimoto.
\newblock Identifying the risks of {LM} agents with an {LM}-emulated sandbox.
\newblock In \emph{Proceedings of ICLR}, 2024.

\bibitem[Sclar et~al.(2024)]{sclar2024quantifying}
M.~Sclar, Y.~Choi, Y.~Tsvetkov, and A.~Suhr.
\newblock Quantifying language models' sensitivity to spurious features in prompt design with relevance-based analysis.
\newblock In \emph{Proceedings of NAACL}, 2024.

\bibitem[Tamkin et~al.(2024)]{tamkin2024evaluating}
A.~Tamkin, K.~Tran, S.~Ganguli, J.~Clark, and D.~Ganguli.
\newblock Evaluating and mitigating discrimination in language model decisions.
\newblock \emph{arXiv preprint arXiv:2312.03689}, 2024.

\bibitem[Willis et~al.(2008)]{willis2008automatic}
A.~Willis, F.~Chantree, A.~de~Roeck, and B.~Nuseibeh.
\newblock Automatic identification of nocuous ambiguity.
\newblock \emph{Research on Language and Computation}, 6:355--374, 2008.

\bibitem[Yao et~al.(2024)]{yao2024tau}
S.~Yao, H.~Chen, J.~Yang, and K.~Narasimhan.
\newblock $\tau^2$-bench: Evaluating multi-turn tool-use agent capabilities.
\newblock \emph{arXiv preprint}, 2024.

\bibitem[Yuan et~al.(2024)]{yuan2024rjudge}
Z.~Yuan, H.~Tang, Y.~Lu, and J.~Liu.
\newblock R-Judge: Benchmarking safety risk awareness for {LLM} agents.
\newblock In \emph{Proceedings of EMNLP}, 2024.

\bibitem[Zhan et~al.(2024)]{zhan2024injecagent}
Q.~Zhan, R.~Liang, S.~Ying, and H.~Sun.
\newblock {InjecAgent}: Benchmarking indirect prompt injections in tool-integrated {LLM} agents.
\newblock In \emph{Findings of ACL}, 2024.

\bibitem[Zhou et~al.(2023)]{zhou2023instruction}
J.~Zhou, T.~Lu, S.~Mishra, S.~Brahma, S.~Basu, Y.~Luan, D.~Tran, and Q.~V. Le.
\newblock Instruction-following evaluation for large language models.
\newblock \emph{arXiv preprint arXiv:2311.07911}, 2023.

\end{thebibliography}

%----------------------------------------------------------------------
\appendix
\section{Perturbation Control Details}
\label{app:perturbation}

Table~\ref{tab:perturbation_detail} reports per-model perturbation sensitivity.
Across three API models and 30 perturbation pairs (15 synonym, 15 reorder), non-ambiguity surface changes increase violation rates from 6.7\% to 13.9\% ($\Delta=7.2\pp$).
Claude shows the lowest perturbation sensitivity ($+4.8\pp$); GPT-5.4 the highest ($+8.3\pp$).

\begin{table}[h]
\centering
\small
\begin{tabular}{@{}lrrr@{}}
\toprule
\textbf{Model} & \textbf{Original} & \textbf{Perturbed} & \textbf{$\Delta$} \\
\midrule
GPT-5.4 & 7.8\% & 16.1\% & $+8.3$ \\
GPT-4.1 & 6.3\% & 12.4\% & $+6.1$ \\
Claude & 4.2\% & 9.0\% & $+4.8$ \\
\midrule
\textbf{Pooled} & \textbf{6.7\%} & \textbf{13.9\%} & $\mathbf{+7.2}$ \\
\bottomrule
\end{tabular}
\caption{Per-model perturbation sensitivity (178 episodes). $\Delta$ in percentage points.}
\label{tab:perturbation_detail}
\end{table}

\section{Reliability Analysis}
\label{app:reliability}

\begin{table}[h]
\centering
\small
\begin{tabular}{@{}llrcl@{}}
\toprule
\textbf{Layer} & \textbf{Type} & \textbf{ICC} & \textbf{95\% CI} & \textbf{Interp.} \\
\midrule
Ling & Scopal & 0.952 & $\CI{0.912}{0.978}$ & Excellent \\
Ling & Lexical & 0.938 & $\CI{0.892}{0.968}$ & Excellent \\
Ling & Coref. & 0.924 & $\CI{0.871}{0.960}$ & Excellent \\
Spec & Auth.\ scope & 0.834 & $\CI{0.748}{0.904}$ & Good \\
Spec & Cond.\ prec. & 0.821 & $\CI{0.732}{0.893}$ & Good \\
Spec & Incomp. & 0.792 & $\CI{0.694}{0.873}$ & Good \\
\midrule
& \textbf{Overall} & \textbf{0.882} & $\CI{0.842}{0.915}$ & \textbf{Good} \\
\bottomrule
\end{tabular}
\caption{Per-type ICC from 30 clauses $\times$ 3 replicates $\times$ 2 conditions $\times$ 2 models (355 episodes).}
\label{tab:icc}
\end{table}

\section{Preliminary Ambiguity Classifier}
\label{app:classifier}

As a proof-of-concept for automated ambiguity detection, we fine-tuned DeBERTa-v3-large on our 300 clause pairs (240 train / 60 test, 80/20 split, 5 epochs, lr=$2\times10^{-5}$).

\begin{table}[h]
\centering
\small
\begin{tabular}{@{}lrrrr@{}}
\toprule
\textbf{Type} & \textbf{$N$} & \textbf{$F_1$} & \textbf{Prec.} & \textbf{Rec.} \\
\midrule
Incompleteness & 10 & 0.86 & 0.88 & 0.84 \\
Auth.\ scope & 10 & 0.82 & 0.84 & 0.80 \\
Lexical & 10 & 0.79 & 0.82 & 0.76 \\
Cond.\ prec. & 10 & 0.76 & 0.78 & 0.74 \\
Scopal & 10 & 0.73 & 0.75 & 0.71 \\
Coreferential & 10 & 0.68 & 0.72 & 0.65 \\
\midrule
\textbf{Overall} & \textbf{60} & \textbf{0.78} & \textbf{0.81} & \textbf{0.75} \\
\bottomrule
\end{tabular}
\caption{DeBERTa-v3-large PoC classifier results (AUC$=0.85$). Per-type $F_1$ correlates strongly with disambiguation token delta ($r=0.91$), suggesting the classifier relies partly on surface-level information asymmetry.}
\label{tab:classifier}
\end{table}

The classifier achieves $F_1=0.78$ (precision$=0.81$, recall$=0.75$, AUC$=0.85$), demonstrating feasibility of automated ambiguity detection even with limited training data.
Per-type performance correlates with disambiguation token delta ($r=0.91$): types requiring more tokens to disambiguate (incompleteness, authorization scope) are easier to detect, while coreferential ambiguity ($F_1=0.68$), which differs by only a few tokens between variants, is hardest.
This suggests that better classifiers will need semantic representations beyond surface features.

\section{Per-Type Bootstrap Stability}
\label{app:bootstrap}

To assess robustness of the type hierarchy (C2), we performed 10{,}000 bootstrap resamples of per-type violation rates.
Incompleteness ranked first in 97.8\% of resamples; coreferential ranked last in 89.4\%.
Mid-tier type rankings were unstable: no mid-tier type held the same rank in more than 62\% of resamples, confirming that fine-grained ordering within the mid tier is unreliable.
The three-tier structure (high / mid / low) was reproduced in 86.2\% of resamples.

\end{document}
